\section{Cel budowania systemu i zakres systemu}

\subsection{Cele główne}

\begin{itemize}[leftmargin=*]
    \item Dopasowanie sposobu podania materiału do stylu uczenia się ucznia (VAK)
    \item Udostępnienie nauczycielowi edytora kursów z materiałami oznaczonymi formą V, A lub K
    \item Aktualizacja modelu ucznia na podstawie wyników quizów i proponowanie zmiany stylu
    \item Wyjaśnialność: każda decyzja systemu ma widoczne uzasadnienie (,,dlaczego to widzisz'')
\end{itemize}

\subsection{Cele szczegółowe}

\begin{itemize}[leftmargin=*]
    \item Wyznaczenie stylu ucznia z ankiety VAK jako punktu startowego, z możliwością ręcznej zmiany
    \item Zachowanie dostępu do wszystkich materiałów -- materiały innych stylów są zwinięte, nie ukryte
    \item Dobór pytań quizu do materiałów, z których uczeń korzystał
    \item Ocenianie quizów wyłącznie po stronie serwera, odporne na wielokrotne podejścia
    \item Udostępnianie kursów uczniom przez katalog kursów opublikowanych oraz przez klasy nauczyciela
    \item Prosty i czytelny interfejs, przystosowany do młodszych użytkowników
\end{itemize}

\subsection{Zakres systemu}

\begin{itemize}[leftmargin=*]
    \item Rejestracja z wyborem roli (uczeń lub nauczyciel), logowanie, sesja
    \item Administracja kont nauczycieli i uczniów (tworzenie, edycja, blokada)
    \item Kursy i lekcje: tworzenie, kolejność, publikacja, archiwizacja
    \item Edytor lekcji: bloki tekstu (formatowanie, tabele, wzory) i załączniki z formami V/A/K
    \item Quiz lekcji z pytaniami wspólnymi i pytaniami o materiały danej formy
    \item Ankieta VAK, profil ucznia (styl, wagi, historia), propozycja zmiany stylu
    \item Klasy nauczyciela: uczniowie i przypisane kursy
\end{itemize}